\scanpage{82}
Dans le cas ``arithmétique'', on voit donc que l'ext. (13) est scindée sss $\unsure{\gamma_i|q_i}=0$ i.e. si l'on peut trouver, globalement sur $S$, une uniformisante sur $X$ le long de $\unsure{q|g}_i(S)$; et cette condition (tjrs satisfaite loc.\inserted{(\unsure{Zar})} sur $S$) est en tous cas tjrs suffisante (cas arith. ou pas !)

Bien plus, parmi les scindages, il y a ceux provenant du choix d'une \deleted{section}\inserted{base} de $\unsure{J_i/J_i^2|I_i/I_i^2}$ qui sont privilégiés. L'indétermination du choix d'une telle base est dans $G_m(S)=H^0(S,G_m)$ --- celle du choix d'une section de (13) est dans
\[
(20)\quad H^1(S,T)\simeq\varprojlim H^1(S,\mu_n)
\]
On a ici des suites exactes de Kummer
\[
\deleted{\unsure{(21)|(22)}}\quad 0\longrightarrow\deleted{\unsure{H^0(S)_n}}G_m(S)_n\longrightarrow H^1(S,\mu_n)\longrightarrow{}_n\operatorname{Pic}(S)\longrightarrow1
\]
d'où par passage à la limite
\[
(21)\quad 0-G_m(S)^\wedge\longrightarrow H^1(S,T)\longrightarrow\varprojlim_n{}_n\operatorname{Pic}S\unsure{\longrightarrow?|\longrightarrow0}
\]
Dans le ``cas arithmétique'' le troisième terme est nul, on trouve donc
\[
\deleted{G_m}\quad H^1(S,T)\simeq G_m(S)^\wedge.
\]
\margin{81}
